% =====================================================================
%  Abstract
% =====================================================================
\chapter*{ABSTRACT}
\addcontentsline{toc}{chapter}{ABSTRACT}

Predictive maintenance of rotating machinery depends on the continuous
acquisition of vibration signals whose raw volume, when stored
indiscriminately, becomes an operational and financial liability. At the same
time, the training of analysts and of artificial-intelligence algorithms in
vibration diagnostics requires varied and controlled defect signals, which is
impractical on real assets. This work presents the design and implementation of
\textbf{ARGUS}, a fault simulator for rotating machinery that generates synthetic
vibration signals with a spectral signature parameterized by fault type and
applies selective persistence of readings based on the \(R^3\) criterion
(frequency, amplitude and phase), with a golden rule for rotational-speed
variation and a Discard Parallelepiped to filter redundant readings.

The system comprises a Python/FastAPI backend with signal-generation,
FFT-spectral-processing, discard-engine and PostgreSQL~16 persistence modules,
and a React frontend presenting a time-signal panel, an FFT spectrum with
normalized orders \(N\) and rotation in Hz, processing and discard-rate
indicators, and fault-snapshot registration. The catalog covers thirteen fault
types with spectral signatures derived from the vibration-diagnostics
literature, including unbalance, angular and parallel misalignment, looseness,
rotor rub, oil-film instabilities and rolling-element bearing faults (BPFO,
BPFI, BSF and FTF).

Validation was carried out through automated tests (63 backend and 12 frontend
tests) and an end-to-end browser trial, with verification of the spectral
signatures produced by the API. The results indicate that the simulator
reproduces the expected signatures and that the discard engine performs
selective persistence according to the defined decision hierarchy, enabling the
creation of training datasets and root-cause-analysis research without exposing
real assets.

\vspace{0.6cm}
\noindent\textbf{Keywords:} predictive maintenance; vibration analysis; FFT;
synthetic signal generation; selective persistence; condition monitoring.
